% Part: turing-machines
% Chapter: undecidability
% Section: trakhtenbrot

\documentclass[../../../include/open-logic-section]{subfiles}

\begin{document}

\olfileid{tur}{und}{tra}
\olsection{مبرهنة تراختنبروت}

\begin{explain}
  سبق في \olref[rep]{sec} تعريف ال!!{sentence}s $!T(\OLId{latin-looped}{M},\OLId{latin-ordinary}{w})$
  و~$!E(\OLId{latin-looped}{M},\OLId{latin-ordinary}{w})$ للآلة~$\OLId{latin-looped}{M}$ وسلسلة الدخل~$\OLId{latin-ordinary}{w}$. ثم أثبتنا في
  \olref[ver]{lem:valid-if-halt} و\olref[ver]{lem:halt-if-valid}
  أن الصدق المنطقي لـ$!T(\OLId{latin-looped}{M},\OLId{latin-ordinary}{w})\lif !E(\OLId{latin-looped}{M},\OLId{latin-ordinary}{w})$ يكافئ توقف~$\OLId{latin-looped}{M}$
  في نهاية المطاف عند الدخل~$\OLId{latin-ordinary}{w}$. فكان امتناع تقرير مسألة
  التوقف مستلزمًا لامتناع تقرير الصدق المنطقي وقابلية الإرضاء
  لل!!{sentence}s من الرتبة الأولى
  (\cref{tur:und:uns:thm:decision-prob,tur:und:uns:cor:undecidable-sat}).

  إلا أن التعريف المتقدم للصدق المنطقي وقابلية الإرضاء يتناول
  !!{structure}s كيفما كانت، منتهية أو غير منتهية. فهل يهون
  التقرير إذا قصرنا النظر على كون !!a{sentence} قابلة للإرضاء
  في !!{structure} منتهية، أو صادقة في جميع
  ال!!{structure}s المنتهية؟ لا؛ وسنعدل حجة استحالة حل مسألة
  القرار ليظهر لنا ذلك.

  ولنرجع أولًا إلى برهان \olref[ver]{lem:halt-if-valid}.
  فقد أنشأنا فيه نموذجًا~$\Struct M$ لـ~$!T(\OLId{latin-looped}{M},\OLId{latin-ordinary}{w})$ يطابق
  عمل الآلة~$\OLId{latin-looped}{M}$ عند الدخل~$\OLId{latin-ordinary}{w}$، وكان مجاله~$\Nat$ غير المنتهي.
  فإن كانت $\OLId{latin-looped}{M}$ تتوقف عند الدخل~$\OLId{latin-ordinary}{w}$، أمكن إنشاء نموذج منته
  $\Struct M'$ بالطريقة نفسها. ولتتوقف $\OLId{latin-looped}{M}$ عند الدخل~$\OLId{latin-ordinary}{w}$
  بعد~$\OLId{latin-ordinary}{k}$ خطوات. نجعل المجال~$\Domain{\OLId{latin-looped}{M}'}$ هو
  $\{\OLMathNumeral{0},\dots,\OLId{latin-ordinary}{n}\}$، حيث $\OLId{latin-ordinary}{n}$ أكبر العددين:~$\OLId{latin-ordinary}{k}+\OLMathNumeral{1}$ وطول~$\OLId{latin-ordinary}{w}$،
  ونضع
  \[
    \Assign{\prime}{\OLId{latin-looped}{M}'}(\OLId{latin-ordinary}{x}) =
    \begin{cases}
      \OLId{latin-ordinary}{x} + \OLMathNumeral{1} &\text{إذا كان $x < n$}\\
      \OLId{latin-ordinary}{n} &\text{خلاف ذلك،}
    \end{cases}
  \]
  ونضع $\tuple{\OLId{latin-ordinary}{x},\OLId{latin-ordinary}{y}}\in\Assign{<}{\OLId{latin-looped}{M}'}$ إذا وفقط إذا كان
  $\OLId{latin-ordinary}{x}<\OLId{latin-ordinary}{y}$ أو $\OLId{latin-ordinary}{x}=\OLId{latin-ordinary}{y}=\OLId{latin-ordinary}{n}$. وأما سائر تعريف $\Struct{M'}$ فكتعريف
  $\Struct{M}$. فيلزم من بناء~$\Struct{M'}$، على نحو ما في
  \olref[ver]{lem:halt-if-valid}، أن $\Sat{\OLId{latin-looped}{M}'}{!T(\OLId{latin-looped}{M},\OLId{latin-ordinary}{w})}$.
  ومن فرض توقف $\OLId{latin-looped}{M}$ عند الدخل~$\OLId{latin-ordinary}{w}$ يلزم $\Sat{\OLId{latin-looped}{M}'}{!E(\OLId{latin-looped}{M},\OLId{latin-ordinary}{w})}$.
  فـ$\Struct{M'}$ نموذج منته لـ$!T(\OLId{latin-looped}{M},\OLId{latin-ordinary}{w})\land !E(\OLId{latin-looped}{M},\OLId{latin-ordinary}{w})$.
  وتنبه إلى إبدال $\lif$ بـ~$\land$.

  وبذلك ثبتت إحدى الجهتين: توقف $\OLId{latin-looped}{M}$ عند الدخل~$\OLId{latin-ordinary}{w}$ يقتضي
  وجود نموذج منته لـ$!T(\OLId{latin-looped}{M},\OLId{latin-ordinary}{w})\land !E(\OLId{latin-looped}{M},\OLId{latin-ordinary}{w})$. لكن عكسها
  باطل؛ فقد لا تتوقف آلة عند دخل~$\OLId{latin-ordinary}{w}$، مع وجود نموذج منته
  لـ$!T(\OLId{latin-looped}{M},\OLId{latin-ordinary}{w})\land !E(\OLId{latin-looped}{M},\OLId{latin-ordinary}{w})$. ومثال ذلك الآلة~$\OLId{latin-looped}{M}$ ذات
  الحالة الوحيدة $\OLId{latin-ordinary}{q}_\OLMathNumeral{0}$ والتعليمة
  $\delta(\OLId{latin-ordinary}{q}_\OLMathNumeral{0},\TMblank)=\tuple{\OLId{latin-ordinary}{q}_\OLMathNumeral{0},\TMblank,\TMstay}$.
  فإنها تبدأ بالدخل الخالي~$\OLId{latin-ordinary}{w}=\emptyseq$، وتدور أبدًا
  من غير كتابة جديدة ولا حركة للرأس؛ فتهيئاتها كلها واحدة:
  حالة واحدة، وموضع رأس واحد، ومحتوى شريط واحد. ومع ذلك يمكن
  تعريف !!{structure} منتهية~$\Struct{M''}$ تحقق
  $!T(\OLId{latin-looped}{M},\emptyseq)\land !E(\OLId{latin-looped}{M},\emptyseq)$، كما في التمرين.
  وسنعالج هذا بتغيير $!T(\OLId{latin-looped}{M},\OLId{latin-ordinary}{w})$ بحيث نمنع مثل هذه
  ال!!{structure}s.

  \begin{prob}
    خذ $\OLId{latin-looped}{M}$ آلة لا حالة لها سوى $\OLId{latin-ordinary}{q}_\OLMathNumeral{0}$، ولا تعليمة سوى
    $\delta(\OLId{latin-ordinary}{q}_\OLMathNumeral{0},\TMblank)=\tuple{\OLId{latin-ordinary}{q}_\OLMathNumeral{0},\TMblank,\TMstay}$. واجعل
    $\Domain{\OLId{latin-looped}{M}''}=\{\OLMathNumeral{0},\OLMathNumeral{1},\OLMathNumeral{2}\}$، و
    $\Assign{\prime}{\OLId{latin-looped}{M}''}(\OLMathNumeral{0})=\Assign{\prime}{\OLId{latin-looped}{M}''}(\OLMathNumeral{1})=\OLMathNumeral{1}$، و
    $\Assign{\prime}{\OLId{latin-looped}{M}''}(\OLMathNumeral{2})=\OLMathNumeral{2}$، و
    $\Assign{<}{\OLId{latin-looped}{M}''}=\{\tuple{\OLMathNumeral{0},\OLMathNumeral{1}},\tuple{\OLMathNumeral{1},\OLMathNumeral{1}},\tuple{\OLMathNumeral{2},\OLMathNumeral{2}}\}$.
    عيّن $\Assign{\Obj Q_{\OLId{latin-ordinary}{q}_\OLMathNumeral{0}}}{\OLId{latin-looped}{M}''}$، و$\Assign{\Obj
    S_{\TMblank}}{\OLId{latin-looped}{M}''}$، و$\Assign{\Obj S_{\TMendtape}}{\OLId{latin-looped}{M}''}$ على أن
    تصدق $!T(\OLId{latin-looped}{M},\emptyseq)$ و$!E(\OLId{latin-looped}{M},\emptyseq)$، وأقم الحجة على صدقهما.
    إرشاد: إن $\delta(\OLId{latin-ordinary}{q}_\OLMathNumeral{0},\TMendtape)$ غير معرّفة. واطلب أن تصدق
    \begin{align*}
    & \Obj Q_{\OLId{latin-ordinary}{q}_\OLMathNumeral{0}}(\num{\OLMathNumeral{1}}, \num{\OLId{latin-ordinary}{n}}) \land \Obj S_{\TMendtape}(\num{\OLMathNumeral{0}}, \num{\OLId{latin-ordinary}{n}})
      \land
      \lforall[\OLId{latin-ordinary}{x}][(\num{\OLMathNumeral{0}} < \OLId{latin-ordinary}{x} \lif \Obj S_\TMblank(\OLId{latin-ordinary}{x}, \num{\OLId{latin-ordinary}{n}}))] \text{\quad لكل $n \in \Nat$}\\
    & \lexists[\OLId{latin-ordinary}{y}][(\Obj Q_{\OLId{latin-ordinary}{q}_\OLMathNumeral{0}}(\num{\OLMathNumeral{0}}, \OLId{latin-ordinary}{y}) \land \Obj S_{\TMendtape}(\num{\OLMathNumeral{0}}, \OLId{latin-ordinary}{y}))]
    \end{align*}
    معًا في~$\Struct{M''}$.
  \end{prob}
\end{explain}

فلنأخذ ال!!{sentence}s الواصفة لعمل الآلة~$\OLId{latin-looped}{M}$ على الدخل
$\OLId{latin-ordinary}{w}=\sigma_{\OLId{latin-ordinary}{i}_\OLMathNumeral{1}}\dots\sigma_{\OLId{latin-ordinary}{i}_\OLId{latin-ordinary}{k}}$، بالترتيب الآتي:
\begin{enumerate}
\item بديهيات الأعداد والعلاقة~$<$، كما في تعريف~$!T(\OLId{latin-looped}{M},\OLId{latin-ordinary}{w})$
في \olref[rep]{sec}.
\item بديهيات تهيئة الدخل نفسها التي في تعريف~$!T(\OLId{latin-looped}{M},\OLId{latin-ordinary}{w})$.
\item بديهيات الانتقال بين التهيئتين المتعاقبتين:

نبقي معنى $!A(\OLId{latin-ordinary}{x},\OLId{latin-ordinary}{y})$ على ما تقدم، ونضع
\[
  !B(\OLId{latin-ordinary}{y}) \ident \lforall[\OLId{latin-ordinary}{x}][(\OLId{latin-ordinary}{x} < \OLId{latin-ordinary}{y} \lif \eq/[\OLId{latin-ordinary}{x}][\OLId{latin-ordinary}{y}])].
\]
\begin{enumerate}
\item \ollabel{rep-right} نقابل كل تعليمة $\delta(\OLId{latin-ordinary}{q}_\OLId{latin-ordinary}{i},\sigma)=
  \tuple{\OLId{latin-ordinary}{q}_\OLId{latin-ordinary}{j},\sigma',\TMright}$ بال!!{sentence}:
\begin{align*}
& \lforall[\OLId{latin-ordinary}{x}][\lforall[\OLId{latin-ordinary}{y}][(
   (\Obj Q_{\OLId{latin-ordinary}{q}_\OLId{latin-ordinary}{i}}(\OLId{latin-ordinary}{x}, \OLId{latin-ordinary}{y}) \land \Obj S_{\sigma}(\OLId{latin-ordinary}{x}, \OLId{latin-ordinary}{y})) \lif {}]] \\
&\qquad   (\Obj Q_{\OLId{latin-ordinary}{q}_\OLId{latin-ordinary}{j}}(\OLId{latin-ordinary}{x}', \OLId{latin-ordinary}{y}') \land \Obj S_{\sigma'}(\OLId{latin-ordinary}{x}, \OLId{latin-ordinary}{y}') \land
!A(\OLId{latin-ordinary}{x}, \OLId{latin-ordinary}{y}) \land !B(\OLId{latin-ordinary}{y}')))
\end{align*}
\item \ollabel{rep-left} نقابل كل تعليمة $\delta(\OLId{latin-ordinary}{q}_\OLId{latin-ordinary}{i},\sigma)=
  \tuple{\OLId{latin-ordinary}{q}_\OLId{latin-ordinary}{j},\sigma',\TMleft}$ بال!!{sentence}
\begin{align*}
& \lforall[\OLId{latin-ordinary}{x}][\lforall[\OLId{latin-ordinary}{y}][
    ((\Obj Q_{\OLId{latin-ordinary}{q}_\OLId{latin-ordinary}{i}}(\OLId{latin-ordinary}{x}', \OLId{latin-ordinary}{y}) \land \Obj S_{\sigma}(\OLId{latin-ordinary}{x}', \OLId{latin-ordinary}{y})) \lif {}]]\\
& \qquad   (\Obj Q_{\OLId{latin-ordinary}{q}_\OLId{latin-ordinary}{j}}(\OLId{latin-ordinary}{x}, \OLId{latin-ordinary}{y}') \land \Obj S_{\sigma'}(\OLId{latin-ordinary}{x}', \OLId{latin-ordinary}{y}') \land
!A(\OLId{latin-ordinary}{x}', \OLId{latin-ordinary}{y}) \land !B(\OLId{latin-ordinary}{y}'))) \land {}\\
& \lforall[\OLId{latin-ordinary}{y}][((\Obj Q_{\OLId{latin-ordinary}{q}_\OLId{latin-ordinary}{i}}(\num{\OLMathNumeral{0}}, \OLId{latin-ordinary}{y}) \land \Obj S_{\sigma}(\num{\OLMathNumeral{0}},
    \OLId{latin-ordinary}{y})) \lif {}]\\
& \qquad (\Obj Q_{\OLId{latin-ordinary}{q}_\OLId{latin-ordinary}{j}}(\num{\OLMathNumeral{0}}, \OLId{latin-ordinary}{y}') \land \Obj S_{\sigma'}(\num{\OLMathNumeral{0}},
  \OLId{latin-ordinary}{y}') \land !A(\num{\OLMathNumeral{0}}, \OLId{latin-ordinary}{y}) \land !B(\OLId{latin-ordinary}{y}')))
\end{align*}
\item \ollabel{rep-stay} نقابل كل تعليمة $\delta(\OLId{latin-ordinary}{q}_\OLId{latin-ordinary}{i},\sigma)=
  \tuple{\OLId{latin-ordinary}{q}_\OLId{latin-ordinary}{j},\sigma',\TMstay}$ بال!!{sentence}:
\begin{align*}
& \lforall[\OLId{latin-ordinary}{x}][\lforall[\OLId{latin-ordinary}{y}][(
   (\Obj Q_{\OLId{latin-ordinary}{q}_\OLId{latin-ordinary}{i}}(\OLId{latin-ordinary}{x}, \OLId{latin-ordinary}{y}) \land \Obj S_{\sigma}(\OLId{latin-ordinary}{x}, \OLId{latin-ordinary}{y})) \lif {}]] \\
&\qquad (\Obj Q_{\OLId{latin-ordinary}{q}_\OLId{latin-ordinary}{j}}(\OLId{latin-ordinary}{x}, \OLId{latin-ordinary}{y}') \land \Obj S_{\sigma'}(\OLId{latin-ordinary}{x}, \OLId{latin-ordinary}{y}') \land
!A(\OLId{latin-ordinary}{x}, \OLId{latin-ordinary}{y}) \land !B(\OLId{latin-ordinary}{y}')))
\end{align*}
\end{enumerate}
وكما ترى بال!!{sentence}s التي تصف انتقالات~$\OLId{latin-looped}{M}$ مماثلة
لل!!{sentence} المقابلة في~$!T(\OLId{latin-looped}{M},\OLId{latin-ordinary}{w})$، إلا أننا نضيف $!B(\OLId{latin-ordinary}{y}')$ في
النهاية. وتضمن $!B(\OLId{latin-ordinary}{y}')$ أن يكون العدد $\OLId{latin-ordinary}{y}'$، رقم التهيئة «التالية»،
مختلفًا عن جميع الأعداد السابقة $\Obj 0$، و$\Obj 0'$، و\dots.
% \item Sentences that express consistency conditions:
% \begin{enumerate}
%   \item\ollabel{rep-con-symb}%
%   !!^a{sentence} that says that no square can have more than one
%   symbol on it at any given time:
%   \[\lforall[x][\lforall[y][(\Obj S_{\sigma}(x, y) \lif \lnot \Obj
%   S_{\sigma'}(x, y))]],\]
%   for every pair $\sigma \neq \sigma'$ of tape symbols of~$T$.
%   \item\ollabel{rep-con-state}%
%   !!^a{sentence} that says that $M$ cannot be in more than one state
%   at any given time:
%   \[\lforall[x][\lforall[y][(\Obj Q_{q}(x, y) \lif
%   \lforall[z][\lnot \Obj
%   Q_{q'}(z, y))]],\]
%   for every pair $q \neq q'$ of states of~$T$.
%   \item\ollabel{rep-con-tape}%
%   !!^a{sentence} that says that $M$ cannot be on more than one tape
%   square at any given time:
%   \[\lforall[x][\lforall[y][(\Obj Q_{q}(x, y) \lif
%   \lforall[z][\eq/[z][y]\lif \lnot \Obj
%   Q_{q'}(z, y))]],\]
%   for every pair $q$, $q'$ of states of~$T$.
% \end{enumerate}
\end{enumerate}
ولنسمّ $!T'(\OLId{latin-looped}{M},\OLId{latin-ordinary}{w})$ اقتران ال!!{sentence}s المتقدمة جميعًا،
الخاصة بالآلة~$\OLId{latin-looped}{M}$ والدخل~$\OLId{latin-ordinary}{w}$.

\begin{lem}\ollabel{lem:halts-sat}
  توقف $\OLId{latin-looped}{M}$ عند الدخل~$\OLId{latin-ordinary}{w}$ يستلزم أن يكون لـ
  $!T'(\OLId{latin-looped}{M},\OLId{latin-ordinary}{w})\land !E(\OLId{latin-looped}{M},\OLId{latin-ordinary}{w})$ نموذج منته.
\end{lem}

\begin{proof}
  ننشئ $\Struct{M'}$ على نحو ما في برهان
  \olref[ver]{lem:halt-if-valid}، غير أننا نجعل
  \begin{align*}
    \Domain{\OLId{latin-looped}{M}'} & = \{\OLMathNumeral{0}, \dots, \OLId{latin-ordinary}{n}\},\\
    \Assign{\prime}{\OLId{latin-looped}{M}'}(\OLId{latin-ordinary}{x}) & =
    \begin{cases}
      \OLId{latin-ordinary}{x} + \OLMathNumeral{1} &\text{إذا كان $x < n$}\\
      \OLId{latin-ordinary}{n} &\text{خلاف ذلك،}
    \end{cases} \\
    \tuple{\OLId{latin-ordinary}{x},\OLId{latin-ordinary}{y}} \in \Assign{<}{\OLId{latin-looped}{M}'} &\text{إذا وفقط إذا كان $x < y$ أو $x = y = n$،}
  \end{align*}
  وفي ذلك $\OLId{latin-ordinary}{n}=\max(\OLId{latin-ordinary}{k}+\OLMathNumeral{1},\len{\OLId{latin-ordinary}{w}})$، و$\OLId{latin-ordinary}{k}$ أقل عدد تتوقف $\OLId{latin-looped}{M}$
  بعده، أي بعد~$\OLId{latin-ordinary}{k}$ خطوات من ابتدائها بالدخل~$\OLId{latin-ordinary}{w}$. ويبقى،
  على سبيل التمرين، إثبات
  $\Sat{\OLId{latin-looped}{M}'}{!T'(\OLId{latin-looped}{M},\OLId{latin-ordinary}{w})\land !E(\OLId{latin-looped}{M},\OLId{latin-ordinary}{w})}$.
\end{proof}

\begin{prob}
  تمم حجة \olref[tur][und][tra]{lem:halts-sat} بإثبات
  $\Sat{\OLId{latin-looped}{M}'}{!T'(\OLId{latin-looped}{M},\OLId{latin-ordinary}{w})\land !E(\OLId{latin-looped}{M},\OLId{latin-ordinary}{w})}$.
\end{prob}

\begin{lem}\ollabel{lem:sat-halts}
  وجود نموذج منته لـ$!T'(\OLId{latin-looped}{M},\OLId{latin-ordinary}{w})\land !E(\OLId{latin-looped}{M},\OLId{latin-ordinary}{w})$ يقتضي أن
  تتوقف $\OLId{latin-looped}{M}$ عند الدخل~$\OLId{latin-ordinary}{w}$.
\end{lem}

\begin{proof}
  نثبت العكس النقيض. لتكن $\OLId{latin-looped}{M}$ غير متوقفة عند~$\OLId{latin-ordinary}{w}$.
  فإن لم يكن لـ$!T'(\OLId{latin-looped}{M},\OLId{latin-ordinary}{w})\land !E(\OLId{latin-looped}{M},\OLId{latin-ordinary}{w})$ نموذج أصلًا،
  حصل المطلوب. وإلا فليكن $\Struct{M}$ نموذجًا
  لـ~$!T'(\OLId{latin-looped}{M},\OLId{latin-ordinary}{w})\land !E(\OLId{latin-looped}{M},\OLId{latin-ordinary}{w})$؛ ونثبت أنه غير منته.

  % تصحيح مصدر (C273-01): لا تلزم B(0) من البديهيات المعروضة؛ اقتصر استعمالها على n>0.
  يمكن، على مثال حجة \olref[ver]{lem:config}، إثبات أن
  $\OLId{latin-looped}{M}$ إذا لم تتوقف عند الدخل~$\OLId{latin-ordinary}{w}$ بعد~$\OLId{latin-ordinary}{n}$ خطوات، كان
  $!T'(\OLId{latin-looped}{M},\OLId{latin-ordinary}{w})\Entails !C(\OLId{latin-looped}{M},\OLId{latin-ordinary}{w},\OLId{latin-ordinary}{n})$؛ وإذا كان $\OLId{latin-ordinary}{n}>\OLMathNumeral{0}$ لزم أيضًا
  $!T'(\OLId{latin-looped}{M},\OLId{latin-ordinary}{w})\Entails !B(\num{\OLId{latin-ordinary}{n}})$.
  ومن عدم توقف $\OLId{latin-looped}{M}$ عند~$\OLId{latin-ordinary}{w}$ نستنتج
  $!T'(\OLId{latin-looped}{M},\OLId{latin-ordinary}{w})\Entails !C(\OLId{latin-looped}{M},\OLId{latin-ordinary}{w},\OLId{latin-ordinary}{n})\land !B(\num{\OLId{latin-ordinary}{n}})$
  لكل~$\OLId{latin-ordinary}{n}>\OLMathNumeral{0}$. ثم إن \olref[rep]{prop:mlessk} تعطي
  $!T'(\OLId{latin-looped}{M},\OLId{latin-ordinary}{w})\Entails\num{\OLId{latin-ordinary}{k}}<\num{\OLId{latin-ordinary}{n}}$ لكل~$\OLId{latin-ordinary}{k}<\OLId{latin-ordinary}{n}$، ولدينا
  $!B(\num{\OLId{latin-ordinary}{n}})\Entails\num{\OLId{latin-ordinary}{k}}<\num{\OLId{latin-ordinary}{n}}\lif
  \eq/[\num{\OLId{latin-ordinary}{k}}][\num{\OLId{latin-ordinary}{n}}]$. فيلزم
  $\Sat{\OLId{latin-looped}{M}}{\eq/[\num{\OLId{latin-ordinary}{k}}][\num{\OLId{latin-ordinary}{n}}]}$ لكل~$\OLId{latin-ordinary}{k}<\OLId{latin-ordinary}{n}$.
  فهذه الحدود~$\num{\OLId{latin-ordinary}{k}}$، وهي غير متناهية العدد، لها في
  $\Struct{M}$ قيم متباينة. ولا يسعها إلا مجال $\Domain{\OLId{latin-looped}{M}}$
  غير منته؛ فلا يكون $\Struct{M}$ نموذجًا منتهيًا
  لـ$!T'(\OLId{latin-looped}{M},\OLId{latin-ordinary}{w})\land !E(\OLId{latin-looped}{M},\OLId{latin-ordinary}{w})$.
\end{proof}

\begin{prob}
  تمم حجة \olref[tur][und][tra]{lem:sat-halts} بإقامة البرهان على أنه إذا لم
  تتوقف~$\OLId{latin-looped}{M}$، حين تبدأ عند الدخل~$\OLId{latin-ordinary}{w}$، بعد~$\OLId{latin-ordinary}{n}>\OLMathNumeral{0}$ خطوة، فإن
  $!T'(\OLId{latin-looped}{M},\OLId{latin-ordinary}{w})\Entails !B(\num{\OLId{latin-ordinary}{n}})$.
\end{prob}

\begin{thm}[مبرهنة تراختنبروت]
  \ollabel{thm:trakhtenbrodt}
  لا يمكن تقرير وجود نموذج منته ل!!{sentence} اعتباطية من
  الرتبة الأولى، أي تقرير قابليتها للإرضاء على نحو منته.
\end{thm}

\begin{proof}
  لو وجدت آلة تورنغ~$\OLId{latin-looped}{F}$ تقرر قابلية الإرضاء المنتهية،
  لأخذنا أي آلة~$\OLId{latin-looped}{M}$ ودخل~$\OLId{latin-ordinary}{w}$، وحسبنا الجملة
  $!T'(\OLId{latin-looped}{M},\OLId{latin-ordinary}{w})\land !E(\OLId{latin-looped}{M},\OLId{latin-ordinary}{w})$، ثم استعملنا~$\OLId{latin-looped}{F}$ في تقرير
  وجود نموذج منته لها. وبحسب
  \cref{tur:und:tra:lem:halts-sat,tur:und:tra:lem:sat-halts}،
  يوجد لها نموذج منته إذا وفقط إذا توقفت $\OLId{latin-looped}{M}$ عند
  الدخل~$\OLId{latin-ordinary}{w}$. فتكون $\OLId{latin-looped}{F}$ سبيلًا إلى حل مسألة التوقف،
  وقد ثبت امتناع حلها.
\end{proof}

\begin{cor}\ollabel{cor:fproof-incomp}%
  % تصحيح مصدر (C273-02): أُظهر شرط الفعالية اللازم لحجة القابلية للحساب.
  يمتنع وجود نظام !!{derivation} فعّال سليم وتام للصدق المنطقي
  \emph{المنتهي}، تكون اشتقاقاته قابلة للتعداد الفعّال (أو يكون فحص
  براهينه المنتهية قابلًا للقرار)؛ أي نظام تتحقق فيه $\Proves !B$ إذا وفقط
  إذا كان $\Sat{\OLId{latin-looped}{M}}{!B}$ لكل !!{structure} منتهية~$\Struct{M}$.
\end{cor}

\begin{proof}
  برهانه متروك للتمرين.
\end{proof}

\begin{prob}
  برهن \olref[tur][und][tra]{cor:fproof-incomp}. استعمل أن
  $!B$ تصدق في كل !!{structure} منتهية إذا وفقط إذا كانت
  $\lnot !B$ غير قابلة للإرضاء على نحو منته. ثم بيّن أن
  قابلية الإرضاء المنتهية شبه قابلة للقرار بالمعنى الوارد في
  \olref[tur][und][uns]{thm:valid-ce}. واستنتج من ذلك أن
  وجود نظام !!a{derivation} فعّال للصدق المنطقي المنتهي كان
  سيجعل قابلية الإرضاء المنتهية قابلة للقرار.
\end{prob}

\end{document}
